\section{Complete-source persistence and the original endpoint}
\label{sec:v76-persistence}
Apply the physical construction to $W=W^\varepsilon$ without
altering the guard. For a fixed choice of the physical roof field
at each exact event, define
\begin{align}
 Q_r^\sigma(t)&=\operatorname*{ess\,sup}_{0<h<r}R_h^\sigma(t),
 &Q_r(t)&=\min\{Q_r^-(t),Q_r^+(t)\},\label{eq:v76-Q}\\
 \mathcal Q_{\varepsilon,M}(r)&=
 \limsup_{m\to\infty}\sup_{R,n,k}m^2
                 \esssup_{t\in\mathcal C_{m,R,n,k}(M)}Q_r(t).
 \label{eq:v76-ordered-Q}
\end{align}
The fibers at translated roofs are on the full roof line. They are
not cut off at the central target. All word and exact-source
contributions are included before any supremum or collision limit.
The field is part of the specified transport, but an auxiliary cap
partition is not.

\begin{theorem}[Persistence of the complete positive source]
\label{thm:v76-persistence}
For almost every $t$ and every fixed $r>0$,
\begin{equation}\label{eq:v76-persistence}
       [b^\varepsilon(t)-2A_r b^\varepsilon(t)]_+\le Q_r(t).
\end{equation}
This holds without a count-uniform exit, drift, or variation bound.
\end{theorem}
\begin{proof}
For either sign, Theorem~\ref{thm:v76-subcoupling} gives
$b^\varepsilon(t)\le b^\varepsilon(t+\sigma h)+Q_r^\sigma(t)$
for almost every $h\in(0,r)$. Integrate in $h$. The one-sided
average of a nonnegative density is at most twice its centered
average. Therefore
$b^\varepsilon(t)\le2A_r b^\varepsilon(t)+Q_r^\sigma(t)$
for each sign, and taking their minimum proves the assertion.
The common almost-everywhere versions follow from the joint
coarea formulation; no supremum over a null exceptional set is used.
\end{proof}

The defect can be bounded directly by physical sources. Put
\[
 \mathcal X_r^\sigma(t)=\int_{\{L=t\}\cap\Omega}
 W^\varepsilon(x)\one_{\{\text{lifetime in direction }\sigma<r\}}
                                      \dd\mu_t(x).
\]
An exit exactly at $r$ is immaterial to an essential supremum over
$h<r$. Monotonicity of the stopped domains and
Proposition~\ref{prop:v76-drift} give
\begin{equation}\label{eq:v76-physical-budget}
 Q_r(t)\le\min_{\sigma\in\{-1,1\}}
 \left\{\mathcal X_r^\sigma(t)
            +\int_0^rG_\sigma(t+\sigma s)\dd s\right\}.
\end{equation}
The expression is permitted to be infinite. It is not declared
small on the basis of a coordinate cover. In a Morse neighborhood
one may use the sharper direct defect \eqref{eq:v76-morse-defect}
instead of this drift majorant.

\begin{theorem}[A quantitative closure implication on the original record]
\label{thm:v76-closure}
Assume Input statement~\ref{input:v76-analytic}. Suppose, in addition,
that the complete stopped physical transports satisfy, for some
$A_M<\infty$, $q\ge0$, $\beta>0$ and fixed $\varepsilon_0,r_0>0$,
\begin{equation}\label{eq:v76-new-premise}
 \mathcal Q_{\varepsilon,M}(r)\le A_M\varepsilon^{-q}r^\beta
 \qquad(0<\varepsilon<\varepsilon_0,\ 0<r<r_0).
\end{equation}
Then the original exact-label pointwise law \eqref{eq:v76-target}
holds. More precisely, for every sufficiently large fixed $B$ and
fixed $r<r_0$,
\begin{equation}\label{eq:v76-error-bound}
 \mathcal E_M\le C'_MB^{-1/192}
                         +A_M\varepsilon(B)^{-q}r^\beta.
\end{equation}
The bound \eqref{eq:v76-new-premise} is an additional dynamical
premise, not a conclusion of the optical or local Morse theorems.
\end{theorem}
\begin{proof}
Combine \eqref{eq:v76-remainder} with
Theorem~\ref{thm:v76-persistence}. Positivity gives
\[
 |m^2p-\mathcal L|\le 2m^2A_r b^{\varepsilon(B)}
                         +m^2Q_r+|e_B|.
\]
Fix $B$, then $\varepsilon(B)$, then $r$, and the auxiliary band
in \eqref{eq:v76-average}, before taking the collision limsup.
Enlarge that auxiliary band only afterwards. The result is
\[
 \mathcal E_M\le 2C\varepsilon(B)^{1/16}
                 +\mathcal Q_{\varepsilon(B),M}(r)+C_MB^{-1/192}.
\]
Since $\varepsilon(B)^{1/16}=A_0^{1/16}B^{-1/192}$,
\eqref{eq:v76-new-premise} proves \eqref{eq:v76-error-bound}.
Choose $r(B)=\varepsilon(B)^{(q+a)/\beta}$ for any $a>0$,
reducing it by a fixed constant if necessary. It is fixed before
the collision limit. Finally let $B\to\infty$. Both terms vanish.
No band, roof width, critical collar, or mesh is chosen as a
function of $m$ inside a spectral estimate.
\end{proof}

\subsection{Downstream statements require scalar closure}
For precision, the retained path input has the following form.
At each exact label and roof there is a finite positive path
remainder $\mathsf B_B(t)$ with mass
$\mathsf B_B(t)(1)=m^2b^{\varepsilon(B)}(t)$. If $\mathsf P_m(t)$
is the normalized path-numerator density and $\mathsf L_m(t)$
its complete arithmetic reference measure, then
\[
 \mathsf P_m-\mathsf L_m=\mathsf B_B+\mathsf e_B,
 \qquad\limsup_m\sup_{R,n,k}\esssup_{\mathcal C(M)}
          \|\mathsf e_B\|_{\mathrm{BL}^*}\le C_MB^{-1/192}.
\]
The bounded-Lipschitz test class includes the constant one. These
are measures for the same scalar disintegration, not separately
selected density versions for different tests.

\begin{corollary}[The same-roof normalization implication]
\label{cor:v76-path}
If \eqref{eq:v76-new-premise} and the preceding path input hold,
the path-numerator error vanishes in the same ordered limits.
Conditional path ratios may then be formed on central arithmetic
classes whose scalar reference is bounded below by a fixed $d>0$.
Arithmetic zero classes are never divided by. Any actual-return
clock-transfer or forward-likelihood conclusion additionally retains
its original hypotheses; none becomes unconditional from a
single-word optical estimate.
\end{corollary}
\begin{proof}
The dual bounded-Lipschitz norm of a positive measure is its mass.
The proof of Theorem~\ref{thm:v76-closure} bounds that mass after
the same limits and therefore bounds the path error. On a class
with reference at least $d$, scalar convergence makes the actual
normalized density at least $d/2$ for large $m$. Subtract the two
ratios, use this lower bound and the convergent numerator and
denominator errors. This proves the conditional implication on
that class only. Clock transfer and likelihood are not proved by
this algebra and require their stated retained inputs.
\end{proof}
